\chapter{课程介绍}\label{preface}
\sourcepage{1}\idx{Langlands纲领}
\textbf{课程：Langlands纲领导论。}\par
\textbf{授课教师：Yao YUAN。}

\textbf{关于本课程。}本课程旨在对Langlands纲领作一个初步介绍。我们主要沿用E. Frenkel的讲义（\bref{mainref1}{1}的第1--5节），同时补充更多细节和预备知识，使没有深厚数论背景的学生也能学习这门课程。

另一方面，尽管这一领域既深刻又广阔，也请不要对教师通过本课程所能讲授内容的深度或广度抱有过高的期待。虽然能力有限，我们仍会尽力而为。

课程共23讲，大致安排如下：
\begin{itemize}
\item 数论中的Langlands纲领。（9讲）
\begin{enumerate}
\item 数域上的类域论。（第\hyperref[L1]{1}--\hyperref[L2]{2}讲）
\item Langlands对应的一般表述。（第\hyperref[L3]{3}--\hyperref[L5]{5}讲）
\item 模形式与Hecke代数。（第\hyperref[L6]{6}--\hyperref[L8]{8}讲）
\item 从数域到函数域。（第\hyperref[L9]{9}讲）
\end{enumerate}
\item 几何Langlands纲领。（14讲）
\begin{enumerate}
\item Langlands对应的几何重述。（第\hyperref[L10]{10}--\hyperref[L12]{12}讲）
\item 反常层与$D$-模。（第\hyperref[L13]{13}--\hyperref[L14]{14}讲）
\item 几何阿贝尔类域论。（第\hyperref[L15]{15}--\hyperref[L17]{17}讲）
\item 其他约化群的Langlands对应。（第\hyperref[L18]{18}--\hyperref[L20]{20}讲）
\item 几何Satake同构。（第\hyperref[L21]{21}--\hyperref[L23]{23}讲）
\end{enumerate}
\end{itemize}

主要参考文献如下。每一讲的末尾还会列出与该讲特别相关的参考文献。
\begin{references}
\refitem{mainref1}{1} E. Frenkel, \textit{Lectures on the Langlands Program and Conformal Field Theory}.\par (\url{https://math.berkeley.edu/~frenkel/houches.pdf})
\sourcepage{2}
\refitem{mainref2}{2} S. Gelbart, \textit{An elementary introduction to the Langlands program}, Bull. Amer. Math. Soc. 10 (1984) 177--219.
\refitem{mainref3}{3} J. Bernstein and S. Gelbart, \textit{An introduction to the Langlands Program}, Birkh\"auser, 2003.
\refitem{mainref4}{4} A. W. Knapp, \textit{Introduction to the Langlands program}, in Representation theory and automorphic forms (Edinburgh, 1996), pp. 245--302, Proc. Symp. Pure Math. 61, AMS, 1997.
\end{references}
